\clearpage
\VerdierPhysicalPageAnchor{0213}%
\begingroup\small\setlength{\parskip}{0pt}\setlength{\emergencystretch}{1.5em}
\setlength{\abovedisplayskip}{4pt plus 2pt minus 2pt}%
\setlength{\belowdisplayskip}{4pt plus 2pt minus 2pt}%

\VerdierSemanticAnchor{ch3d-3-2-5-second-diagram}
\[
\begin{tikzcd}[column sep=4.5em]
X & Z^\bullet
  \arrow[l,"{u(S^n(X,X'))}"']
  \arrow[r,"{v(S^n(X,X'))}"]
& X'[n]
\end{tikzcd}
\quad,
\]
したがって次の図を得る：

\VerdierSemanticAnchor{ch3d-3-2-5-composite-diagram}
\begin{center}
\resizebox{0.98\textwidth}{!}{%
$\begin{tikzcd}[ampersand replacement=\&,column sep=3.0em]
X
\& Z^\bullet
  \arrow[l,"{u(S^n)}"']
  \arrow[r,"{v(S^n)}"]
\& X'[n]
\& {Z'}^\bullet[n]
  \arrow[l,"{u(S^m)[n]}"']
  \arrow[r,"{v(S^m)[n]}"]
\& X''[m+n]
\end{tikzcd}$% 
}
\end{center}
（記法を簡単にするため、明らかな略記を用いた。）したがって、
$D(\mathcal A)$ におけるシフトの定義（第 II 章，2.2.6）から
次の等式を得る：

\VerdierSemanticAnchor{ch3d-3-2-5-1}
\begin{equation}
\begin{aligned}
&\bigl(\delta^m(S^m(X',X''))\bigr)[n]\,
  \delta^n(S^n(X,X'))\\
&\qquad={}
[v(S^m)[n]][u(S^m)[n]]^{-1}[v(S^n)][u(S^n)]^{-1}.
\end{aligned}
\tag{3.2.5.1}
\end{equation}

\VerdierSemanticAnchor{ch3d-3-2-5-2}
ここで次のようにおく：
\begin{equation}
\begin{aligned}
Z''^\bullet={}&\cdots\longrightarrow0\longrightarrow X''
\xrightarrow{(-1)^n u'_m}Z'_{m-1}
\xrightarrow{(-1)^n u'_{m-1}}\cdots\\
&\cdots\xrightarrow{(-1)^n u'_1}Z'_0
\xrightarrow{u_nu'_0}Z_{n-1}
\xrightarrow{u_{n-1}}\cdots\xrightarrow{u_1}Z_0
\longrightarrow0\longrightarrow0\longrightarrow\cdots.
\end{aligned}
\tag{3.2.5.2}
\end{equation}
（$Z_i$ は次数 $-i$ にある $Z''^\bullet$ の成分である）、

\VerdierSemanticAnchor{ch3d-3-2-5-spliced-extension}
\[
\begin{aligned}
S^{m+n}(X,X'')={}&0\longrightarrow X''
\xrightarrow{(-1)^n u'_m}Z'_{m-1}
\xrightarrow{(-1)^n u'_{m-1}}\cdots\\
&\cdots\xrightarrow{(-1)^n u'_1}Z'_0
\xrightarrow{u_nu'_0}Z_{n-1}
\xrightarrow{u_{n-1}}\cdots\xrightarrow{u_1}Z_0
\xrightarrow{u_0}X\longrightarrow0.
\end{aligned}
\]
ここで複体の射を二つ定める：

\VerdierSemanticAnchor{ch3d-3-2-5-3}
\begin{equation}
\begin{aligned}
\rho\colon Z''^\bullet\longrightarrow Z^\bullet,
&\qquad
\begin{cases}
\rho^p=\operatorname{id},&p\geq -n+1,\\
\rho^{-n}=u'_0,&\\
\rho^p=0,&p<-n;
\end{cases}\\[3pt]
\psi\colon Z''^\bullet\longrightarrow {Z'}^\bullet[n],
&\qquad
\begin{cases}
\psi^p=0,&p\geq -n+1,\\
\psi^p=\operatorname{id},&p\leq -n.
\end{cases}
\end{aligned}
\tag{3.2.5.3}
\end{equation}
次の性質が直ちに確かめられる：

\VerdierSemanticAnchor{ch3d-3-2-5-a}
\noindent a) 射
$\rho\colon Z''^\bullet\longrightarrow Z^\bullet$ (3.2.5.3) は
擬同型である。
\endgroup
